% !TEX root = ../notes.tex

\documentclass[../notes.tex]{subfiles}

\begin{document}

\section{April 27}
Today, we classify cluster algebras.

\subsection{Cluster Algebras, Again}
Last time, we constructed a cluster algebra attached to a skew-symmetrizable matrix. Let's return to this definition, as follows.
\begin{definition}[mutation]
	Fix a positive integer $n$. A \textit{seed} is a pair $(x,B)$ of an $n$-tuple $x=(x_1,\ldots,x_n)$ of algebraically independent variables over $\QQ$, and a \textit{mutation matrix} $B$ is an $n\times n$ skew-symmetric matrix. Given this data, we define the $k$th \textit{mutation} $\mu_k(x,B)$ to be the seed given by
	\[\mu_k(x)_i\coloneqq\begin{cases}
		x_i & \text{if }i=k, \\
		\frac{\displaystyle\prod_{b_{jk}>0}x_j^{\left|b_{jk}\right|}+\displaystyle\prod_{b_{jk}<0}x_j^{\left|b_{jk}\right|}}{\displaystyle x_k} & \text{if }i\ne k,
	\end{cases}\]
	together with the mutation matrix given by the matrix $B'$ with coefficients
	\[B'_{ij}=\begin{cases}
		-b_{ij} & \text{if }i=k\text{ or }j=k, \\
		b_{ij}+\frac{\left|b_{ik}\right|b_{kj}+b_{ik}\left|b_{kj}\right|}2 & \text{if }i\ne k\text{ and }j\ne k.
	\end{cases}\]
\end{definition}
\begin{remark}
	Some of these definitions also work with skew-symmetrizable matrices $B$ (and the formulae have been written to accommodate this), but our seeds in the sequel will mostly come from skew-symmetric matrices $B$.
\end{remark}
\begin{remark}
	By construction, applying the $k$th mutation twice does nothing. This property explains the rule for mutating the mutation matrix. Roughly speaking, the mysterious second equation in the mutation of $B$ should be split up into two cases: if $b_{ik}$ and $b_{kj}$ have the same sign, then the two pieces combine; and if they have different signs, then we have cancellation.
\end{remark}
\begin{definition}[cluster algebra]
	Fix a positive integer $n$, and choose a seed $(x,B)$. Then the associated \textit{cluster algebra} $\mc A$ is the $\ZZ$-subalgebra of $\QQ(x)$, generated by all possible sequence of mutations. The collection of generators is known as the set of \textit{cluster variables}.
\end{definition}
\begin{remark}
	Our cluster algebras $\mc A$ will have $\mc A\cap\QQ=\ZZ$, which follows from \Cref{thm:laurent}.
\end{remark}
It is not at all obvious if the associated cluster algebra is finitely generated. In light of our cluster variables, the correct notion is as follows.
\begin{definition}[finite type]
	A cluster algebra $\mc A$ has \textit{finite type} if and only if there are only finitely many distinct cluster variables.
\end{definition}
Here is how we will test this.
\begin{definition}
	A cluster algebra $\mc A$ is \textit{acyclic} if and only if the associated seed $(x,B)$ (and its mutations) produce directed graphs with no oriented cycles. Here, the directed graph associated to an $n\times n$ skew-symmetric $B$ is the one with vertices $\{v_1,\ldots,v_n\}$ and $b_{ij}$ edges from $v_i$ to $v_j$.
\end{definition}
\begin{remark}
	It turns out that $\mc A$ has finite type if and only if $B$ is acyclic.
\end{remark}
\begin{example}
	The matrix $\begin{bsmallmatrix}
		0 & 1 \\ -1 & 0
	\end{bsmallmatrix}$ induces the graph $1\to 2$. The only mutation of this matrix is its negative, so we see that this cluster algebra is acyclic.
\end{example}
\begin{nex}
	The matrix
	\[\begin{bmatrix}
		0 & 2 & -2 \\ -2 & 0 & 2 \\ 2 & -2 & 0
	\end{bmatrix}\]
	produces the graph
	% https://q.uiver.app/#q=WzAsMyxbMCwwLCIxIl0sWzEsMCwiMiJdLFsxLDEsIjMiXSxbMCwxLCIiLDAseyJvZmZzZXQiOjF9XSxbMCwxLCIiLDIseyJvZmZzZXQiOi0xfV0sWzEsMiwiIiwwLHsib2Zmc2V0IjoxfV0sWzEsMiwiIiwwLHsib2Zmc2V0IjotMX1dLFsyLDAsIiIsMCx7Im9mZnNldCI6MX1dLFsyLDAsIiIsMSx7Im9mZnNldCI6LTF9XV0=&macro_url=https%3A%2F%2Fraw.githubusercontent.com%2FdFoiler%2Fnotes%2Fmaster%2Fnir.tex
	\[\begin{tikzcd}[cramped]
		1 & 2 \\
		& 3
		\arrow[shift right, from=1-1, to=1-2]
		\arrow[shift left, from=1-1, to=1-2]
		\arrow[shift right, from=1-2, to=2-2]
		\arrow[shift left, from=1-2, to=2-2]
		\arrow[shift right, from=2-2, to=1-1]
		\arrow[shift left, from=2-2, to=1-1]
	\end{tikzcd}\]
	which is notably not acyclic. The mutations of this matrix correspond to Vieta jumps. Namely, specializing the seed variables $(x_1,x_2,x_3)\mapsto(1,1,1)$ discovers all solutions to the Markov equation $x^2+y^2+z^2=3xyz$.
\end{nex}

\subsection{Theorems on Cluster Algebras}
Here are the main theorems on cluster algebras. The following is due to Fomin and Zelevinsky.
\begin{theorem}[Laurent phenomenon] \label{thm:laurent}
	Fix a cluster algebra $\mc A$ with initial seed $((x_1,\ldots,x_n),B)$. Then any cluster variable in $\mc A$ is a Laurent polynomial in the variables $(x_1,\ldots,x_n)$. In other words, $\mc A\subseteq\ZZ[x_1,\ldots,x_n,1/x_1,\ldots,1/x_n]$.
\end{theorem}
\begin{proof}[Sketch]
	The key ingredient is the following intersection lemma: if $x$ mutates to $x'$ via the $i$th mutation, then by construction, $x_ix_i'$ can be expanded into a sum of two disjoint monomials. Thus, setting $R\coloneqq\ZZ[x_1^\pm,\ldots,x_{i-1}^\pm,x_{i+1}^\pm,\ldots,x_n^\pm]$, we find
	\[R\left[x_i^\pm\right]\cap R\left[(x_i')^\pm\right]=R[x_i,x_i'].\]
	For example, the inclusion from the right to the left follows from the expansion of $x_ix_i'$. The moral is that having a cluster variable $z$ which is a Laurent polynomial in the variables $(x_1,\ldots,x_n)$ and also a Laurent polynomial in all adjacent mutated seeds, then $z\in\ZZ[x_1,\ldots,x_n,x_1',\ldots,x_n']$, so $z$ is Laurent.
	
	One can then do an induction in the ``distance'' of a cluster variable from the initial seed. Namely, at distance zero, there is nothing to do. At distance one, this follows from the construction. At distance two, one has a cluster variable which is Laurent in distance one, and one can check by hand that the resulting polynomials are also Laurent in all the other mutations in distance two (and hence polynomial), which again uses the given mutation formula explicitly. After this, the induction kicks in: for example, in distance three, one can use the distance-two version to show that a given cluster variable is a polynomial in cluster variables that are distance zero and one.
\end{proof}
\begin{remark}
	Fomin called this proof ``elementary but technical but elementary.''
\end{remark}
\begin{remark}
	Here is an application: if we set $x_1=\cdots=x_n=1$, then all the cluster variables will be sent to integers. This remains true as long as the $x_i$s are sent to elements of a ring which are units.
\end{remark}
The following is due to many people, starting with Lee--Schiffler and ending with Gross--Hacking--Keel--Kontesevich. It uses scattering diagrams from mirror symmetry, which we do not have time for, so we will say nothing about its proof.
\begin{theorem}[Positivity phenomenon]
	Fix a cluster algebra $\mc A$ with initial seed $((x_1,\ldots,x_n),B)$. Then any cluster variable in $\mc A$ is a Laurent polynomial in the variables $(x_1,\ldots,x_n)$ with nonnegative coefficients.
\end{theorem}
\begin{remark}
	The recursion can directly show that these polynomials take positive values on positive points.
\end{remark}
Our last result is classification.
\begin{definition}[Cartan companion]
	Fix a skew-symmetric matrix $B$. Then we define the \textit{Cartan companion} $C$ to be the matrix $C$ with coefficients
	\[c_{ij}\coloneqq\begin{cases}
		2 & \text{if }i=j, \\
		-\left|b_{ij}\right| & \text{if }i\ne j.
	\end{cases}\]
\end{definition}
\begin{example}
	The Cartan companion of the matrix $\begin{bsmallmatrix}
		0 & 1 \\ -1 & 0
	\end{bsmallmatrix}$ is the matrix $\begin{bsmallmatrix}
		2 & -1 \\ -1 & 2
	\end{bsmallmatrix}$.
\end{example}
\begin{theorem}
	Fix a cluster algebra $\mc A$ with initial seed $((x_1,\ldots,x_n),B)$. Then $\mc A$ has finite type if and only if $B$ and all of its iterated mutations has positive-definite Cartan companion.
\end{theorem}
\begin{proof}[Sketch]
	In rank $2$, our $B$ looks like $\begin{bsmallmatrix}
		0 & b \\ -c & 0
	\end{bsmallmatrix}$, and we are interested in what happens to $(x_1,x_2)$ when we alternatively apply the two mutations $\mu_1$ and $\mu_2$. This gives a sequence $\{x_1,x_2,x_3,\ldots\}$ of cluster variables. They satisfy
	\[x_{m-1}x_{m+1}=\begin{bsmallmatrix}
		x_m^b+1 & \text{if }m\text{ is odd}, \\
		x_m^c+1 & \text{if }m\text{ is even}.
	\end{bsmallmatrix}\]
	One finds that the sequence $\{x_n\}$ is periodic (which is equivalent to finite type in this case because there are not so many mutations) if and only if $\left|bc\right|\le3$: the backward direction is a direct calculation, and the forward direction follows by some degree argument.

	We now move on to higher rank. Let's start with the forward direction. We will show that if the Cartan companion avoids one of the finite Dynkin types fails to have finite type. Well, one should start by checking that $B$ does not admit one of the rank-$2$ situations. It turns out that all affine Dynkin types (away from the finite type Dynkin types) and other infinite types have one of these infinite-type rank-$2$ subalgebras.

	Lastly, we must show that $B$ having finite Dynkin type implies that $\mc A$ has finite type, which is some direct calculation. Let $\Phi$ be the associated finite root system, and we choose a polarization with simple roots $\Pi=\{\alpha_1,\ldots,\alpha_n\}$. Now, for any cluster variable $z$, \Cref{thm:laurent} tells us that
	\[z=\frac{P(x_1,\ldots,x_n)}{x_1^{d_1}\cdots x_n^{d_n}},\]
	where $P$ is not divisible by any $x_i$ with $d_i\ne0$. We can now define a map which sends the cluster variables to $\op{span}\Phi$ which sends $x_i\mapsto-\alpha_i$ for each $i$ and in general sends $z$ to $\sum_id_i\alpha_i$. It turns out that $z$ gets sent to a positive root by this map, and it turns out that this map is injective, so there are only finitely many cluster variables. In total, we find that the cluster variables are in bijection with $\Phi^+\sqcup-\Pi$, so we being finite type follows.
\end{proof}
\begin{example}
	One can check that \Cref{ex:clusterl-alg-a3} is the cluster algebra from type $A_3$.
\end{example}
\begin{nex}
	The seed $\left(x_1,x_2,\begin{bsmallmatrix}
		0 & -3 \\ 3 & 0
	\end{bsmallmatrix}\right)$ has all mutations of the mutation matrix with Cartan companion $\begin{bsmallmatrix}
		2 & 3 \\ 3 & 2
	\end{bsmallmatrix}$, which fails to be positive-definite. Indeed, this matrix admits a negative eigenvalue.
\end{nex}
\begin{remark}
	It turns out that the Cartan companions being positive-definite can be checked from the initial seed. This follows from the property of its mutation.
\end{remark}

\end{document}